\documentclass[10pt,a4paper]{amsart}
\usepackage[margin=19mm]{geometry}
\usepackage{amsmath,amssymb,mathtools}
\usepackage[scr=rsfs,cal=euler]{mathalfa}
\usepackage{xfrac}
\usepackage[unicode,hidelinks,bookmarks=false]{hyperref}
\newcommand{\ie}{i.e.\ }
\newcommand{\cf}{cf.\ }
\newcommand{\resp}{respectively\ }
\newcommand{\Subsection}{\subsection}
\providecommand{\nobreakdash}{\nobreak\hbox{-}}
\setlength{\emergencystretch}{2em}
\multlinegap0pt
\usepackage{tikz}
\usepackage{amssymb}
\usepackage{mathtools}
\newtheorem{theorem}{Theorem}
\title{Convergence Rate of Projected Subgradient Method with Time-varying Step-sizes}
\author{Zhihan Zhu, Yanhao Zhang, Yong Xia}
\date{Source: arXiv:2402.10221v1 (CC BY 4.0)}
\begin{document}
\maketitle

\noindent\emph{Source and licence.} Convergence Rate of Projected Subgradient Method with Time-varying Step-sizes. Version arXiv:2402.10221v1 (CC BY 4.0), distributed by arXiv under the Creative Commons Attribution 4.0 International licence, http://creativecommons.org/licenses/by/4.0/, as recorded in the arXiv licence metadata for this exact version and shown on the abstract page https://arxiv.org/abs/2402.10221v1. Full licence notice: \texttt{licenses/arxiv-2402.10221-CC-BY-4.0.txt}.\par\smallskip

\noindent\textit{Editorial scope.} Counted unit: the article's theorem Convergence (source0.tex lines 188--194). The article's complete original proof of it is delivered with it (source0.tex lines 195--235, one \texttt{proof} environment of 41 lines). No mathematics is written, repaired or re-derived: the statement and the proof are the article's own lines, in the article's own notation, and no retained line is substituted or re-flowed.  No further source-local context is needed: every cross-reference of every retained line resolves inside the two delivered spans.  Nothing was omitted from any retained span, so the package carries no omission record.  No retained line contains a citation, so the wrapper carries no bibliography and no bibliography entry.  No retained line contains a figure environment, an image-inclusion command or drawing code, so no image is delivered and the package declares no asset dependency.  Editorial formatting only, in addition: an AMS \texttt{amsart} preamble that restates the article's own declarations and macros used by the retained lines, namely the environments \texttt{theorem} on separate counters, together with the standard packages those lines need.  The retained environments print in this document as Theorem 1 in order of appearance, whereas the article prints the same items with the numbering of its own full document.  The complete native source of the article as distributed in the arXiv e-print for this exact version is bundled unchanged as \texttt{sources/154635/source0.tex}.
\begin{theorem}\label{Convergence}
For any fixed $k\ge -1$,
PSG with a positive and non-increasing step-size sequence $\{\eta_s\}$ satisfies
	\begin{equation}
		f\left(\frac{\sum_{s=1}^t \frac{1}{\eta_s^k}x_s}{\sum_{s=1}^t \frac{1}{\eta_s^k}}\right)-f(x^*)\le \frac{\frac{R^2}{\eta_t^{k+1}}+ L^2 \sum_{s=1}^t\frac{1}{\eta_s^{k-1}}}{2\sum_{s=1}^t\frac{1}{\eta_s^k}}. \label{bound}
	\end{equation}
\end{theorem}	

\begin{proof}
	According to the definition of subgradient, we have
	\begin{eqnarray}
			f(x_s)-f(x^*)&\le& g_s^T(x_s-x^*) \nonumber \\
			&=& \frac{1}{\eta_s}(x_s-y_{s+1})^T(x_s-x^*) \nonumber\\
			&=& \frac{1}{2\eta_s}(\Vert x_s-y_{s+1} \Vert^2+\Vert x_s-x^* \Vert^2-\Vert y_{s+1}-x^* \Vert^2)\label{eq1}\\
			&=& \frac{1}{2\eta_s}(\Vert x_s-x^* \Vert^2-\Vert y_{s+1}-x^* \Vert^2)+\frac{\eta_s}{2}\Vert g_s \Vert^2 \nonumber\\
&\le& \frac{1}{2\eta_s}(\Vert x_s-x^* \Vert^2-\Vert x_{s+1}-x^* \Vert^2)+\frac{\eta_s}{2}L^2,\label{ineq1}
	\end{eqnarray}
where \eqref{eq1} follows from the elementary identity $2a^Tb=\Vert a\Vert^2+\Vert b\Vert^2-\Vert a-b\Vert^2$, and \eqref{ineq1} holds since $\Vert g_s \Vert \le L$ and
	\begin{equation*}
		\Vert y_{s+1}-x^* \Vert^2\ge \Vert x_{s+1}-x^* \Vert^2,
	\end{equation*}
which is implied by  the projection theorem.
%	Since $f$ is $L$-Lipschitz on $\mathcal{X}$, $\Vert g_s \Vert \le L$, we have
%	\begin{equation*}
%		f(x_s)-f(x^*)\le \frac{1}{2\eta_s}(\Vert x_s-x^* \Vert^2-\Vert x_{s+1}-x^* \Vert^2)+\frac{\eta_s}{2}L^2.
%	\end{equation*}
%Consequently, for any $k\ge -1$,
%	\begin{equation*}
%		\frac{1}{\eta_s^k}(f(x_s)-f(x^*))\le \frac{1}{2\eta_s^{k+1}}(\Vert x_s-x^* \Vert^2-\Vert x_{s+1}-x^* \Vert^2)+\frac{1}{2\eta_s^{k-1}}L^2.
%	\end{equation*}
%	Summing up the inequality over $s$, we get

Consequently, we have
	\begin{eqnarray}
&&\left(\sum_{s=1}^t\frac{1}{\eta_s^k}\right)\left(		f\left(\frac{\sum_{s=1}^t \frac{1}{\eta_s^k}x_s}{\sum_{s=1}^t \frac{1}{\eta_s^k}}\right)-f(x^*)\right) \nonumber \\
&\le&	\sum_{s=1}^t\frac{1}{\eta_s^k}(f(x_s)-f(x^*))  ~~~~({\rm since~}f{\rm~is~convex})\nonumber\\
&\le& \sum_{s=1}^t\frac{1}{2\eta_s^{k+1}}(\Vert x_s-x^* \Vert^2-\Vert x_{s+1}-x^* \Vert^2)+\sum_{s=1}^t\frac{1}{2\eta_s^{k-1}}L^2 \label{ineq2}\\
			&=& \frac{1}{2\eta_1^{k+1}}\Vert x_1-x^* \Vert^2+\sum_{s=2}^t(\frac{1}{2\eta_s^{k+1}}-\frac{1}{2\eta_{s-1}^{k+1}})\Vert x_s-x^* \Vert^2 \nonumber\\
&&-\frac{1}{2\eta_{t}^{k+1}}\Vert x_{t+1}-x^* \Vert^2+\sum_{s=1}^t\frac{1}{2\eta_s^{k-1}}L^2\nonumber\\
			&\le& \frac{R^2}{2\eta_1^{k+1}}+\frac{R^2}{2}\sum_{s=2}^t(\frac{1}{\eta_s^{k+1}}-\frac{1}{\eta_{s-1}^{k+1}})+
\sum_{s=1}^t\frac{1}{2\eta_s^{k-1}}L^2\label{ineq3}\\
			&=& \frac{R^2}{2\eta_t^{k+1}}+\sum_{s=1}^t\frac{1}{2\eta_s^{k-1}}L^2,\nonumber
		 \end{eqnarray}
%	Since $f$ is convex on $\mathcal{X}$,
%	\begin{equation*}
%(\sum_{s=1}^t\frac{1}{\eta_s^k})(		f(\frac{\sum_{s=1}^t \frac{1}{\eta_s^k}x_s}{\sum_{s=1}^t \frac{1}{\eta_s^k}})-f(x^*))\le \frac{\frac{R^2}{2\eta_t^{k+1}}+\frac{L^2}{2}\sum_{s=1}^t\frac{1}{\eta_s^{k-1}}}{\sum_{s=1}^t\frac{1}{\eta_s^k}}.
%	\end{equation*}
where \eqref{ineq2} follows from \eqref{ineq1}, and \eqref{ineq3} holds since $1/\eta_s^{k+1}-1/\eta_{s-1}^{k+1}\ge0$ when $k\ge -1$. The proof is complete.  \qed
\end{proof}	

\end{document}
